\documentclass{article}
\usepackage{empheq}
\usepackage{float}
\usepackage{epsfig}
\usepackage{graphicx}
\usepackage{textcomp}
% \usepackctage{epstopdf}
\usepackage{amsmath}
\usepackage{amssymb}
\usepackage[mathscr]{euscript}
\usepackage[margin=1in]{geometry}
\usepackage{framed}
\usepackage{color}
\usepackage{comment}
\linespread{1.5}

\title{Analytical attempt at the 2D problem}
\author{Kevin Bedoya \and Oleg Kogan}
\date{November 2025}

\begin{document}

\maketitle

\section{Full problem recap}
We have two layers - a disk-shaped continuous diffusive layer (DL) with radius $1$, and an advective layer (AL) consisting of discrete filaments or microtubules (MT), with inward-pointing advective velocity.  We will restrict ourselves here to highly symmetric situations, with angularly-symmetric initial conditions, and evenly-spaced advective filaments that meet at the center of the disk.  This invites the use of polar coordinates.
\\
\\
Let $\phi(r,\theta,t)$ be the density on the DL, and $\rho_i(r,t)$ be the density on the $i$-th MT, each located at angle $\theta_i$.  The dynamics of these densities is given by
\begin{eqnarray}
\label{Eq:one} \frac{\partial \phi}{\partial t} &=& \nabla^2 \phi +\left(b \rho_i(r) - a \phi(r,\theta)\right)\sum_{i}\frac{1}{r}\delta(\theta - \theta_i) \\
\label{eq:two} \frac{\partial \rho_i}{\partial t} &=& v \frac{\partial \rho_i}{\partial r} + a\phi(r,\theta_i) - b \rho_i(r)
\end{eqnarray}
(the dependence on $t$ was not explicitly stated to lighten the notation).  This will be derived in the Appendix, but for now, it's important to mention that parameter $a$ has dimensions of $distance/time$, and parameter $b$ has dimensions of $1/time$.  Also, here $D$ is set to $1$ (which amounts to a change of time units), and the radius of the disk is also set to $1$ (change of length units).  The sum of $\delta$-functions says that DL receives and loses mass to AL only at the angles and which the MTs are laid down.  

The boundary conditions (BC)s are absorbing, meaning $\phi(r=1) = 0$ and $\rho(r=1) = 0$.  The initial condition (IC) will depend only on the radius (i.e. it is angularly-symmetric), and we will work primarily with a $\delta$-function IC on the DK, $\phi(t=0) = \delta(r)$, and $\rho(t=0) = 0$.

\section{Mapping the problem to a wedge (sector).}
While this formulation is correct, and it is what's used in discrete implementations in a PDE solver on a polar grid, working with $\delta$-functions is a pain.  We do know that for such symmetric IC: (i) the densities on all the MTs will be the same and (ii) the density on the DL will be periodic with period $\frac{2 \pi}{N}$, where $N$ is the number of MTs.  Even more specifically, the MTs divide the DL into wedges (or sectors), and $\phi$ will be periodic from sector to sector. Therefore, we want to reformulate the problem just to a sector.
\\
\\
Consider a sector between $\theta_i$ and $\theta_{i+1}$ (and the sectors adjoint to it.  To extract BCs obeyed just by that sector, we integrate Eq.~(\ref{Eq:one}) through $\theta_i$:
\begin{displaymath}
\int_{\theta_i - \epsilon}^{\theta_i + \epsilon} \frac{\partial \phi}{\partial t} \,d\theta = \int_{\theta_i - \epsilon}^{\theta_i + \epsilon} \left(\nabla^2 \phi +\left(b \rho_i(r) - a \phi(r,\theta)\right)\sum_{i}\frac{1}{r}\delta(\theta - \theta_i)\right) \,d\theta
\end{displaymath}
On the left hand side (LHS), we have
\begin{displaymath}
\int_{\theta_i - \epsilon}^{\theta_i + \epsilon} \frac{\partial \phi}{\partial t} \,d\theta  = \frac{\partial}{\partial t} \int_{\theta_i - \epsilon}^{\theta_i + \epsilon} \phi(r,\theta,t)\,d\theta 
\end{displaymath}
Similar to my thinking (which proved to be correct) in 1D, the density on the DL will not contain $\delta$-function-like terms due to the diffusion operator;  in other words, all singular features in which there's a finite amount of mass concentrated in an infinitesimal amount of space are immediately smoothed out by diffusion.  Then the integral in the last expression will go to zero as $\epsilon \rightarrow 0$.  So the LHS is zero.
\\
\\
Now on to the RHS. 
\begin{eqnarray*}
\int_{\theta_i - \epsilon}^{\theta_i + \epsilon} a \phi(r,\theta) \frac{1}{r}\delta(\theta - \theta_i) \,d\theta &=& \frac{a}{r} \phi(r,\theta_i) \\
\int_{\theta_i - \epsilon}^{\theta_i + \epsilon} b \rho_i(r) \frac{1}{r} \delta(\theta - \theta_i) \,d\theta &=& \frac{b}{r} \rho_i(r)
\end{eqnarray*}
Next, consider 
\begin{displaymath}
\int_{\theta_i - \epsilon}^{\theta_i + \epsilon} \nabla^2 \phi  \,d\theta = \int_{\theta_i - \epsilon}^{\theta_i + \epsilon} \left(\frac{\partial^2 \phi}{\partial r^2} + \frac{1}{r} \frac{\partial \phi}{\partial r} + \frac{1}{r^2}\frac{\partial^2 \phi}{\partial r^2} \right) \,d\theta
\end{displaymath}
The first term is 
\begin{displaymath}
\int_{\theta_i - \epsilon}^{\theta_i + \epsilon} \frac{\partial^2 \phi}{\partial r^2} \,d\theta = \frac{\partial^2}{\partial r^2} \int_{\theta_i - \epsilon}^{\theta_i + \epsilon} \phi(r,\theta,t) \,d\theta,
\end{displaymath}
and the integral goes to zero for the same reason as stated above, so this term is zero.  The same can be said for 
\begin{displaymath}
\int_{\theta_i - \epsilon}^{\theta_i + \epsilon} \frac{1}{r}\frac{\partial \phi}{\partial r} \,d\theta = 0.
\end{displaymath}
Finally,
\begin{displaymath}
\frac{1}{r^2} \int_{\theta_i - \epsilon}^{\theta_i + \epsilon} \frac{\partial^2 \phi}{\partial \theta^2} \,d\theta = \frac{1}{r^2}\left(\left. \frac{\partial \phi}{\partial \theta} \right|_{\theta_i + \epsilon} - \left. \frac{\partial \phi}{\partial \theta} \right|_{\theta_i - \epsilon} \right).
\end{displaymath}
Now, because of the periodicity properties of $\phi$, this can be rewritten entirely in terms of $\phi$ within the wedge, without invoking $\phi$ at a neighbiring wedge.  Thus, the last term can be written as 
\begin{displaymath}
 \frac{1}{r^2}\left(\left. \frac{\partial \phi}{\partial \theta} \right|_{\theta_i + \epsilon} - \left. \frac{\partial \phi}{\partial \theta} \right|_{\theta_{i+1} - \epsilon} \right).
\end{displaymath}
So, the derivative here is evaluated just a little bit inward of the ``left'' side of the wedge, and just a little bit inward of the ``right'' side of the wedge.
\\
\\
All put together, this integration through the angle $\theta_i$ has produced the following relationship:
\begin{equation}
\frac{a}{r} \phi(r,\theta_i)= \frac{1}{r^2}\left(\left. \frac{\partial \phi}{\partial \theta} \right|_L - \left. \frac{\partial \phi}{\partial \theta} \right|_R \right) + \frac{b}{r} \rho(r)
\end{equation}
I'm no longer putting a subscript under $\rho$ to emphasize that density profiles on all MTs are the same.   For the same reason, I called the locations at which the $\theta$-derivatives of $\phi$ are evaluated by letters ``L'' and ``R'' (left and right edges).  It makes it more compact, and avoids notational clutter.  
\\
\\
What we have derived is in fact a new boundary condition for $\phi$ in the wedge.  The other BCs are: (i) $\rho(r=1) = 0$ and periodic BC $\phi(\theta) = \phi\left(\theta+\frac{2\pi}{N}\right)$.  If we commit to only modeling $\phi$ within the wedge, the last BC becomes $\phi(0) = \phi\left(\frac{2\pi}{N}\right)$.   In other words, we placed the left edge of the wedge at $\theta=0$ without loss of generality.
\\
\\
Let's summarize now our restated problem.  
\begin{eqnarray}
\label{eq:one_new} \frac{\partial \phi}{\partial t} &=& \nabla^2 \phi, \\
\label{eq:two_new} \frac{\partial \rho}{\partial t} &=& v \frac{\partial \rho}{\partial r} + a\phi(r,\theta =0) - b \rho(r),
\end{eqnarray}
The BCs are
\begin{eqnarray}
\phi(r=1) &=& 0  \\
\phi(\theta=0) &=& \phi(\theta = \frac{2\pi}{N})  \\
a \phi(r,\theta=0) &=& \frac{1}{r}\left(\left. \frac{\partial \phi}{\partial \theta} \right|_{\theta=0} - \left. \frac{\partial \phi}{\partial \theta} \right|_{\theta=\frac{2\pi}{N}} \right) + b \rho(r) \\
\label{eq:BC4} \rho(r=1) &=& 0
\end{eqnarray}
The first and fourth are the absorbing BC for $\phi$ and $\rho$, the second is the periodic BC for $\phi$, and the third is the new BC.  Also, by symmetry, $-\left. \frac{\partial \phi}{\partial \theta} \right|_{\theta=\frac{2\pi}{N}} = \left. \frac{\partial \phi}{\partial \theta} \right|_{\theta=0}$, which will simplify the third BC.
\\
\\
There's a two-way coupling. The AL density $\rho$ affects the DL density through the BCs on the edges.  On the other hand, the DL density affects AL density directly through the second term on the right.
\\
\\
We are interested in eigenfunction expansion.  In other words, in functions of the form 
\begin{equation}
\left(\begin{array}{c} \phi \\ \rho \end{array}\right) = \left(\begin{array}{c} F(r,\theta) \\ R(r) \end{array}\right) e^{-\sigma t}.
\end{equation}
This is the same as separation of variables into spatial and temporal parts.  When this ansatz is substituted into the equations above, we get the following equation for eigenfunctions $(F,R)$:
\begin{eqnarray}
\label{eq:one_new} -\sigma F &=& \nabla^2 F \\
\label{eq:two_new} -\sigma R &=& v \frac{d R}{d r} + aF(r,\theta=0) - b R(r)
\end{eqnarray}
With BCs:
\begin{eqnarray}
\label{eq:BC1} F(r=1) &=& 0  \\
\label{eq:BC2} F(\theta=0) &=& F(\theta = \frac{2\pi}{N})  \\
\label{eq:BC3} a F(r,\theta=0) &=& \frac{2}{r} \left. \frac{\partial F}{\partial \theta} \right|_{\theta=0}  + b R(r) \\
\label{eq:BC4} R(r=1) &=& 0
\end{eqnarray}

\section{Attempts at solving this problem (newer)}
\noindent Consider Eq.~(\ref{eq:one_new}) along with the first three BCs.  The term $b R(r)$ in the third BC, couples the DL to the AL.  At the same time, Eq.~(\ref{eq:two_new}) is coupled to the DL through the term $aF(r,\theta=0)$.  One exercise that we could perform, is to assume that $R(r)$ is a known function, solve for $F(r,\theta)$ in terms of this function (so, it would be a \emph{functional} of $R(r)$), plug it into the Eq.~(\ref{eq:two_new}), which now becomes a closed equation for $R(r)$.  Our attempts to do this are described below, but they have not been successful (see ``Gameplan'' below).
\\
\\
An alternative approach is reverse - to solve Eq.~(\ref{eq:two_new}) for $R(r)$ as a \emph{functional} of the unknown $F(r,\theta=0)$, and then plug it into Eq.~(\ref{eq:BC3}), which, together with Eq.~(\ref{eq:one_new}) will now be a closed system for $F(r,\theta)$.  This should be easier to do, because Eq.~(\ref{eq:two_new}) is first order, and in one variable.  I will now attempt this approach.
\\
\\
An equation of the type $\frac{dy}{dx} = \alpha y(x) + f(x)$ with $y(1)=0$ has a solution of the form
\begin{displaymath}
y(x) = e^{\alpha x}\int_1^x e^{-\alpha x'} f(x') \,dx'
\end{displaymath}
Applying this to Eq.~(\ref{eq:two_new}) we have
\begin{displaymath}
R(r) = -\frac{a}{v}e^{\frac{b-\sigma}{v} r}\int_1^r e^{-\frac{b-\sigma}{v} r'} F(r',\theta=0) \,dr'
\end{displaymath}
We have extracted everything from Eq.~(\ref{eq:two_new}) and the fourth BC.  We now have a closed problem for $F(r,\theta)$ only.  It obeys the Helmholtz equation $-\sigma F = \nabla^2 F$ with three BCs
\begin{eqnarray}
\label{eq:bc1} F(r=1) &=& 0  \\
\label{eq:bc2} F(\theta=0) &=& F(\theta = \frac{2\pi}{N})  \\
\label{eq:bc3} a F(r,\theta=0) - \frac{2}{r} \left.\frac{\partial F}{\partial \theta} \right|_{\theta=0} &=& -\frac{ab}{v}e^{\frac{b-\sigma}{v} r}\int_1^r e^{-\frac{b-\sigma}{v} r'} F(r',\theta=0) \,dr' 
\end{eqnarray}
We will attempt to solve the Helmholtz equation using the separation of variables $F(r,\theta) = \Phi(r)\Theta(\theta)$.  
We have 
\begin{displaymath}
-\sigma \Phi(r)\Theta(\theta) = \Theta(\theta)\Phi_{rr}(r) + \frac{1}{r}\Phi_{r}(r)\Theta(\theta) + \frac{1}{r^2}\Phi(r)\Theta_{\theta \theta},
\end{displaymath}
or equivalently,
\begin{displaymath}
\sigma r^2 + r^2\frac{\Phi_{rr}}{\Phi} + r \frac{\Phi_r}{\Phi} = -\frac{\Theta_{\theta \theta}}{\Theta}
\end{displaymath}
This says that $\Theta_{\theta \theta} = - \mu^2 \theta$ and $\Phi_{rr} + \frac{1}{r} \Phi_r + \left(\sigma -\frac{\mu^2}{r^2} \right)\Phi  = 0$, where $\mu^2$ is some separation constant.  The $\Theta$-equation has a solution $\Theta = Ae^{i\mu \theta} + Be^{-i\mu \theta}$.  The $\Phi$-equation has a solution $J_{\mu}(\sqrt{\sigma} r)$ (we're not including the Neumann functions $Y$, because those diverge at the origin, and we're seeking a regular solution).  
\\
\\
I will apply the first two boundary conditions to each eigenfunction.  In other words, 
\begin{eqnarray}
\label{eq:bc1} \Phi(r=1) &=& 0  \\
\label{eq:bc2} \Theta(\theta=0) &=& \Theta\left(\theta = \frac{2\pi}{N}\right) 
\end{eqnarray}
Therefore, we will choose $\mu = nN$ to satisfy the secnd of these BCs.  The first BC says that $J_{nN}(\sqrt{\sigma}) = 0$. Let $\kappa_{n,j}$ denote the $j$-th positive root of the Bessel function $J_{nN}$, i.e. $J_{nN}(\kappa_{nN,j}) = 0$ with $j = 1,2,\dots.$
Thus, $\sqrt{\sigma} = \kappa_{nN,j}$.  For each pair $(n,j)$, the corresponding radial basis function is therefore $\Phi_{n,j}(r)
=J_{nN}\left(\kappa_{nN,j}r\right)$.  We note also that $J_n(x) = (-1)^nJ_{-n}(x)$ for integer $n$, and moreover $J_n(ax) = J_n(-ax)$.  Thus, a solution to $-\sigma F = \nabla^2 F$ that satisfy the first two boundary conditions is of the form $e^{i nN\theta}J_{nN}\bigl(\kappa_{nN,j} r\bigr)$, where $n \in\mathbb{Z}$.  
\\
\\
We are left with the third BC:
\begin{displaymath}
\label{eq:bc3} a F(r,\theta=0) - \frac{2}{r} \left.\frac{\partial F}{\partial \theta} \right|_{\theta=0} = -\frac{ab}{v}e^{\frac{b-\kappa_{|n|N, j}^2}{v} r}\int_1^r e^{-\frac{b-\kappa_{|n|N, j}^2}{v} r'} F(r',\theta=0) \,dr' 
\end{displaymath}
We now substitute our eigenfunction $e^{i nN\theta}J_{nN}\bigl(\kappa_{nN,j} r\bigr)$ and get
\begin{equation}
\label{eq:thirdBC_long}
\left(a  - \frac{2inN}{r} \right)J_{|n|N}(\kappa_{|n|N, j} r) + \frac{ab}{v}e^{\frac{b-\kappa_{|n|N, j}^2}{v} r}\int_1^r e^{-\frac{b-\kappa_{|n|N, j}^2}{v} r'} J_{|n|N}(\kappa_{|n|N, j} r') \,dr'  = 0
\end{equation}
And this should apply for each $(n,j)$ term.  My understanding of this is that it imposes the restriction on those $\kappa$s for which this expression evaluates to zero.  In other words, only those $\kappa$s that satisfy this criterion will enter the sum.  However, what bothers me is that the expression is a function of $r$.  So it's really that for some $n$ and $j$, the \emph{function} of $r$ on the left hand side should be zero. \textbf{If} this is possible, i.e. if there exist some select $\{n\}$ and $\{j\}$ for which this is true, then the general solution will be written as 
\begin{equation}
\phi(r,\theta,t) = \sum_{n', j'} C_{n', j'} e^{i n'N\theta}J_{n'N}\bigl(\kappa_{n'N,j'} r\bigr) e^{-\kappa^2_{n'N,j'} t}
\end{equation}
where $n'$ and $j'$ satisfy Eq.~(\ref{eq:thirdBC_long}).  Similarly,
\begin{equation}
R(r) = \sum_{n', j'} D_{n',j'}\left(e^{\frac{b-\kappa_{n'N,j'}^2}{v} r}\int_1^r e^{-\frac{b-\kappa_{n'N,j'}^2}{v} r'} J_{n'N}\bigl(\kappa_{n'N,j'} r'\bigr) \,dr'\right)e^{-\kappa^2_{n'N,j'} t}
\end{equation}
(I absorbed $\frac{a}{v}$ into $D$).  The coefficients $C$ and $D$ will come from initial conditions. 
\newpage
\section{(Older) gameplan}
We exchanged the problem on a disk with $\delta$-functions, to a problem on a wedge, where the effect of the MTs has been rolled into a BC-3 for $\rhi$.  This is still a difficult problem, just expressed without the $\delta$-functions.  
\\
\\
To tackle it, I propose the following gameplan.  We will first consider the following simpler problem.
\begin{equation}
-\sigma F &=& \nabla^2 F 
\end{equation}
with BCs: (i) $F(r=1) = 0$, (ii) $F(\theta=0) = F(\theta = \frac{2\pi}{N})$, (iii) $F(r,\theta=0) = f(r)$.  
Thus, we are looking at uncoupled diffusion equation, which is driven at the boundaries by some \emph{specified} function of radius.  This is like the the simplified version of BC (Eq.~(\ref{eq:BC3})).  Once we understand how to solve this equation, i.e. how to express the solution in terms of a some generic $f(r)$ (which at this point is assumed \emph{given}), we will rewrite it in terms of the \emph{unknown} $R(r)$, and substitute the solution (that will be \emph{in principle} expressed in terms of $R(r)$ back into the $R$-equation with its own BC.  We also need to solve for $\sigma$.  My intuition is that $\sigma$ should be determined by \emph{both} equations, i.e. the solution to the $R$-equation, once the implicit solution for $F$ in terms of the unknown $R$ is substituted into its RHS, is only possible for specific, discrete set of $\sigma$s.
Lastly, we will repeat the whole process with the $\theta$-derivatives included in the third BC.


\section{An attempt to solve the simplified uncoupled problem}

We seek separated solutions of the form
\begin{equation}
F(r,\theta) = R(r)\,\Theta(\theta).
\label{eq:separation-ansatz}
\end{equation}

In polar coordinates, the Laplacian is
\begin{equation}
\Delta F = F_{rr} + \frac{1}{r}F_r + \frac{1}{r^2}F_{\theta\theta}.
\label{eq:laplacian-polar}
\end{equation}

Substituting \eqref{eq:separation-ansatz} into \eqref{eq:laplacian-polar} gives
\begin{equation}
\Delta F
= \Theta(\theta) R''(r)
+ \frac{1}{r}\Theta(\theta) R'(r)
+ \frac{1}{r^2} R(r)\Theta''(\theta).
\end{equation}

Dividing by \(R(r)\Theta(\theta)\) and multiplying by \(r^2\), we obtain
\begin{equation}
r^2 \frac{R''(r)}{R(r)}
+ r \frac{R'(r)}{R(r)}
+ \frac{\Theta''(\theta)}{\Theta(\theta)}.
\label{eq:pre-separation}
\end{equation}

To obtain solutions in terms of Bessel functions, we consider the eigenvalue problem
\begin{equation}
\Delta F = -\sigma F,
\label{eq:eigenvalue-problem}
\end{equation}
which yields
\begin{equation}
r^2 \frac{R''(r)}{R(r)}
+ r \frac{R'(r)}{R(r)}
+ \frac{\Theta''(\theta)}{\Theta(\theta)}
= -\sigma r^2.
\label{eq:separated-equation}
\end{equation}

Separating variables, we introduce a separation constant \(\mu \in \mathbb{R}\) and write
\begin{equation}
r^2 \frac{R''(r)}{R(r)}
+ r \frac{R'(r)}{R(r)}
+ \sigma r^2
= -\frac{\Theta''(\theta)}{\Theta(\theta)}
=: \mu.
\label{eq:separation-constant}
\end{equation}

The angular equation is therefore
\begin{equation}
\Theta''(\theta) + \mu\,\Theta(\theta) = 0.
\label{eq:angular-ode}
\end{equation}

For \(\mu > 0\), writing \(\sqrt{\mu}\) for the positive square root, the general solution is
\begin{equation}
\Theta(\theta)
= A\,e^{-i\sqrt{\mu}\,\theta}
+ B\,e^{i\sqrt{\mu}\,\theta},
\qquad A,B \in \mathbb{R}.
\label{eq:angular-solution}
\end{equation}

The radial equation becomes
\begin{equation}
R''(r) + \frac{1}{r}R'(r)
+ \left(\sigma - \frac{\mu}{r^2}\right) R(r) = 0,
\label{eq:radial-ode}
\end{equation}
which is Bessel’s equation of order \(\nu = \sqrt{\mu}\).
Choosing the solution regular at \(r=0\), we obtain
\begin{equation}
R(r) = D\,J_{\sqrt{\mu}}\!\left(\sqrt{\sigma}\,r\right),
\qquad D \in \mathbb{R}.
\label{eq:radial-solution}
\end{equation}

\subsection*{\textbf{Angular boundary condition}}

We impose the angular periodicity condition
\begin{equation}
F(0,r) = F\!\left(\frac{2\pi}{N}, r\right),
\qquad N \in \mathbb{N}.
\label{eq:angular-bc}
\end{equation}

From the separated form, this requires
\begin{equation}
\Theta(0) = \Theta\!\left(\frac{2\pi}{N}\right).
\end{equation}

Substituting \eqref{eq:angular-solution}, a sufficient condition for nontrivial solutions is
\begin{equation}
e^{\pm i\sqrt{\mu}(2\pi/N)} = 1,
\end{equation}
which holds if
\begin{equation}
\sqrt{\mu} = nN,
\qquad n \in \mathbb{Z},
\end{equation}
so that
\begin{equation}
\mu = (nN)^2.
\label{eq:mu-quantized}
\end{equation}

The angular eigenfunctions are therefore
\begin{equation}
\Theta_n(\theta)
= A_n e^{i nN \theta}
+ B_n e^{-i nN \theta}.
\label{eq:angular-modes}
\end{equation}

\subsection*{\textbf{Radial boundary condition}}

Substituting \(\mu = (nN)^2\) into \eqref{eq:radial-solution}, the radial modes take the form
\begin{equation}
R_n(r)
= D_n\,J_{nN}\!\left(\sqrt{\sigma}\,r\right).
\label{eq:radial-modes}
\end{equation}

We now impose the boundary condition
\begin{equation}
F(\theta,1) = 0.
\label{eq:radial-bc}
\end{equation}

For a fixed angular index \(n\), this implies
\begin{equation}
R_n(1) = 0,
\end{equation}
or equivalently
\begin{equation}
J_{nN}\!\left(\sqrt{\sigma}\right) = 0.
\label{eq:bessel-zero-condition}
\end{equation}

Let \(\kappa_{n,j}\) denote the \(j\)-th positive zero of \(J_{nN}\), i.e.
\[
J_{nN}(\kappa_{n,j}) = 0,
\qquad j = 1,2,\dots
\]
Then
\begin{equation}
\sigma_{n,j} = \kappa_{n,j}^2,
\qquad n \in \mathbb{Z},\; j \in \mathbb{N}.
\label{eq:eigenvalues}
\end{equation}

The corresponding radial eigenfunctions are
\begin{equation}
R_{n,j}(r)
= D_{n,j}\,J_{nN}\!\left(\kappa_{n,j}\,r\right).
\label{eq:final-radial-modes}
\end{equation}



\paragraph{Remark on the role of \(\sigma\).}
The boundary condition at \(r=1\) does not pick out a single value of \(\sigma\); instead, it quantizes \(\sigma\) into a \emph{discrete spectrum} of eigenvalues
\[
\sigma_{n,j} = \kappa_{n,j}^{\,2},
\]
indexed by the angular mode \(n\) and the radial node index \(j\). In other words, \(\sigma\) plays the role of an eigenvalue, and the boundary value problem admits a family of eigenpairs \(\bigl(\sigma_{n,j}, F_{n,j}\bigr)\), where the radial part is built from
\[
J_{nN}\bigl(\kappa_{n,j} r\bigr)
\]
for all allowed orders \(nN\) and all zeros \(\kappa_{n,j}\) of the corresponding Bessel function. The full solution is then obtained as a superposition over all such modes \((n,j)\), rather than being associated with a single root for a fixed Bessel order.
\\
\\
\color{blue}
I find it very strange that $\sigma$s are determined just by the diffusive BC.  
\color{black}

\subsection*{Series representation of the solution}

We now assemble the separated eigenfunctions into a full series solution. For each angular index 
\(n = 0,1,2,\dots\) and radial index \(j = 1,2,\dots\), we have the separated mode
\begin{align}
F_{n,j}(r,\theta)
&=
\bigl(
A_{n,j} \, e^{i nN\theta}
+
B_{n,j} \, e^{-i nN\theta}
\bigr)\,
J_{nN}\bigl(\kappa_{nN,j} r\bigr), \tag{15)}
\end{align}
where \(\kappa_{nN,j}\) denotes the \(j\)-th positive zero of the Bessel function \(J_{nN}\). A general solution
satisfying the homogeneous boundary conditions can then be written as the double series
\begin{align}
F(r,\theta)
&=
\sum_{n=-\infty}^{\infty}
\sum_{j=1}^{\infty}
\bigl(
A_{n,j} \, e^{i nN\theta}
+
B_{n,j} \, e^{-i nN\theta}
\bigr)\,
J_{nN}\bigl(\kappa_{nN,j} r\bigr). \tag{16)}
\end{align}
The coefficients \(A_{n,j}\) and \(B_{n,j}\) are determined by the remaining boundary condition
\(F(\theta=0,r) = f(r)\) via an appropriate expansion of \(f\) in this eigenbasis.

Note that $J_n(x) = (-1)^nJ_{-n}(x)$ for integer $n$. With this in mind, we can write the sum as 
\begin{align}
F(r,\theta)
&=
\sum_{n\in\mathbb{Z}}
\sum_{j=1}^{\infty}
C_{n,j}\,
e^{i nN\theta}\,
J_{|n|N}\bigl(\kappa_{|n|N,j} r\bigr). \tag{17)}
\end{align}
where \(\{C_{n,j}\}_{n\in\mathbb{Z},\, j\in\mathbb{N}}\) are some complex coefficients to be determined.  Here the dependence of the radial part on \(n\) enters only through the order \(|n|N\) of the Bessel
function and its associated zeros \(\kappa_{|n|,j}\). 

\subsection*{Third boundary condition}

We now have to apply the third BC, i.e.\ fit
\[
F(r,\theta=0) = f(r).
\]
Note that \(F\bigl(r,\theta = \frac{2\pi}{N}\bigr) = f(r)\) will be automatically
satisfied by the angular periodicity already built into the eigenfunctions.

\medskip

\noindent
\textbf{Step 1: Start from the full series representation.}

We use the complex exponential form of the separated solution:
\begin{align}
F(r,\theta)
&=
\sum_{n\in\mathbb{Z}}
\sum_{j=1}^{\infty}
C_{n,j}\,
e^{i nN\theta}\,
J_{|n|N}\bigl(\kappa_{|n|N,j} r\bigr), \tag{26)}
\end{align}
where \(\kappa_{|n|N,j}\) is the \(j\)-th positive zero of \(J_{|n|N}\).
Evaluating at \(\theta=0\) gives
\begin{align}
F(r,0)
&=
\sum_{n\in\mathbb{Z}}
\sum_{j=1}^{\infty}
C_{n,j}\,
J_{|n|N}\bigl(\kappa_{|n|N,j} r\bigr)
=
f(r). \tag{27)}
\end{align}

\newpage
\section{Leftovers}
\medskip

\noindent
\textbf{Step 2: Angular orthogonality (projection onto a fixed mode \(n\)).}

To isolate a single angular mode, we multiply \eqref{26)} by
\(e^{-i nN\theta}\) and integrate over one angular period
\(\theta \in [0,2\pi/N]\):
\begin{align}
\int_0^{2\pi/N} F(r,\theta)\,e^{-i nN\theta}\,d\theta
&=
\sum_{m\in\mathbb{Z}}\sum_{j=1}^{\infty}
C_{m,j}\,J_{|m|N}\bigl(\kappa_{|m|N,j} r\bigr)
\int_0^{2\pi/N} e^{i (m-n)N\theta}\,d\theta. \tag{28)}
\end{align}
By orthogonality of the complex exponentials,
\[
\int_0^{2\pi/N} e^{i (m-n)N\theta}\,d\theta
=
\begin{cases}
\dfrac{2\pi}{N}, & m=n,\\[0.4em]
0, & m\neq n,
\end{cases}
\]
so the sum over \(m\) collapses and we obtain
\begin{align}
\int_0^{2\pi/N} F(r,\theta)\,e^{-i nN\theta}\,d\theta
&=
\frac{2\pi}{N}
\sum_{j=1}^{\infty}
C_{n,j}\,J_{|n|N}\bigl(\kappa_{|n|N,j} r\bigr). \tag{29)}
\end{align}
Define the angular projection
\begin{align}
g_n(r)
:=
\frac{N}{2\pi}
\int_0^{2\pi/N} F(r,\theta)\,e^{-i nN\theta}\,d\theta.
\tag{30)}
\end{align}
Then, for each fixed \(n\in\mathbb{Z}\),
\begin{align}
g_n(r)
&=
\sum_{j=1}^{\infty}
C_{n,j}\,J_{|n|N}\bigl(\kappa_{|n|N,j} r\bigr). \tag{31)}
\end{align}

\medskip

\noindent
\textbf{Step 3: Radial orthogonality (projection onto a fixed radial mode \(j\)).}

For a fixed angular index \(n\) and a fixed radial index \(p\in\mathbb{N}\),
multiply \eqref{31)} by \(r\,J_{|n|N}\bigl(\kappa_{|n|N,p} r\bigr)\) and integrate
from \(0\) to \(1\):
\begin{align}
\int_0^1 r\,g_n(r)\,J_{|n|N}\bigl(\kappa_{|n|N,p} r\bigr)\,dr
&=
\sum_{j=1}^{\infty} C_{n,j}
\int_0^1 r\,J_{|n|N}\bigl(\kappa_{|n|N,j} r\bigr)\,
J_{|n|N}\bigl(\kappa_{|n|N,p} r\bigr)\,dr. \tag{32)}
\end{align}
Using the standard Bessel orthogonality relation
\begin{align}
\int_0^1 r\,J_{\nu}\bigl(\kappa_{\nu,j} r\bigr)\,
J_{\nu}\bigl(\kappa_{\nu,p} r\bigr)\,dr
&=
\begin{cases}
0, & j\neq p,\\[0.4em]
\dfrac{1}{2}\,J_{\nu+1}\bigl(\kappa_{\nu,p}\bigr)^2, & j=p,
\end{cases}
\quad \nu = |n|N, \tag{33)}
\end{align}
only the term \(j=p\) survives in \eqref{32)}, giving
\begin{align}
\int_0^1 r\,g_n(r)\,J_{|n|N}\bigl(\kappa_{|n|N,p} r\bigr)\,dr
&=
C_{n,p}\,\frac{1}{2}\,
J_{|n|N+1}\bigl(\kappa_{|n|N,p}\bigr)^2. \tag{34)}
\end{align}
Hence, for each \(n\in\mathbb{Z}\) and \(p\in\mathbb{N}\),
\begin{align}
C_{n,p}
&=
\frac{2}{J_{|n|N+1}\bigl(\kappa_{|n|N,p}\bigr)^2}
\int_0^1 r\,g_n(r)\,J_{|n|N}\bigl(\kappa_{|n|N,p} r\bigr)\,dr. \tag{35)}
\end{align}

\medskip

\noindent
\textbf{Step 4: Expressing the coefficients in terms of the boundary data \(f(r)\).}

The projection \(g_n(r)\) is defined from the full solution \(F(r,\theta)\) in
\eqref{30)}. On the boundary ray \(\theta=0\), we are given
\[
F(r,0) = f(r).
\]
If the boundary data \(f\) is independent of \(\theta\) and we seek a solution
with the same angular symmetry, then only the \(n=0\) mode is present:
\[
g_0(r) = f(r), \qquad g_n(r) \equiv 0 \text{ for } n\neq 0.
\]
In that case, the nonzero coefficients reduce to
\begin{align}
C_{0,p}
&=
\frac{2}{J_{1}\bigl(\kappa_{0,p}\bigr)^2}
\int_0^1 r\,J_{0}\bigl(\kappa_{0,p} r\bigr)\,f(r)\,dr, \tag{36)}
\end{align}
and \(C_{n,p}=0\) for all \(n\neq 0\). More generally, if one is given boundary
data \(f(r,\theta)\) along the full arc \(0\leq\theta\leq 2\pi/N\), the same
procedure with \(g_n(r)\) defined by \eqref{30)} yields the coefficients
\(C_{n,p}\) in terms of the angular Fourier–Bessel projections of \(f\).


\end{document}

